\scanpage{219}

\mspage{93}

et la flèche can.

\[
A_\alpha\xrightarrow{j_\alpha}B
\]

est dans $C$, car les $j_\beta$ le sont et $C$ est stable par limites \inserted{ind.} filtrantes. \deleted{On a \unsure{ainsi}}

On a ainsi achevé la description du système transfini. Pour \deleted{\unsure{tt}} $\alpha$, on a

\[
\mathrm{card}\,\alpha\leq\mathrm{card}\,\unsure{\mathrm{Ss}|\mathrm{Sos}}(B),
\]

puisque l'application $\beta\longmapsto A_\beta$ \deleted{est \unsure{de}} de l'ens $I_\alpha=\{\beta\mid\beta<\alpha\}$ dans $\mathrm{Ss}(B)$ est \uline{str} croissante. Donc il existe un ordinal $\alpha$ tel que $\mathrm{card}\,\alpha\leq\mathrm{card}\,\mathrm{Ss}(B)$, tel que $A_\alpha=B$. \deleted{Pour prouver que $j\in$} Mais alors $i\in\widetilde C_0$, comme limite inductive \inserted{(ou composée)} transfinie de flèches $\in\widetilde C_0$. (Il suffit de la condition $\forall\beta,\ i_{\beta+1,\beta}\in\widetilde C_0$ pour impliquer que tous les morphismes de transition $j_{\beta',\beta}$ sont dans $\widetilde C_0$.)

Revenons à 7.1, \deleted{\unsure{1°}} \inserted{$2^\circ)$}. Il reste à prouver que si $\mathrm{Mon}$ est stable par cochgts. de base, alors on a \deleted{$\mathrm{Mon}$}

\[
\mathrm{Mon}=\widetilde{\mathrm{Mon}_0}.
\]

Mais comme $\mathrm{Mon}$ est \inserted{déjà} stable par composition, et par limites inductives filtrantes \unsure{quelc.|finales|finies}, \deleted{\unsure{et cofacteurs}} \inserted{(donc aussi par facteurs directs)}, on voit par le cor. 2 du th. 1, \S\,1 \inserted{(p.~\unsure{19|18})} que $\mathrm{Mon}=\widetilde{\mathrm{Mon}_0}$. L'énoncé de factorisation est le th 1 du \S\,1.
% archive sheet number: 218
